\section{Biomedical entity normalization and resources}\label{section:background-biomedical-entity-normalization-and-resources}

Papers use different terminology for the same concepts, so a normalization process is necessary for the pipeline to compare findings across papers. This is the task of biomedical entity normalization: it maps surface forms that occur in papers' text to a controlled vocabulary. This thesis combines scispaCy biomedical models~\parencite{neumann2019scispacy} with the Unified Medical Language System (UMLS)~\parencite{bodenreider2004umls} to assign concept identifiers (CUIs) to subject and outcome entities, as described in Section~\ref{section:Knowledge_Extraction}.

The assigned CUIs provide the comparability gate with a stable way to determine whether two findings refer to the same entity. If both findings have a CUI, the gate compares those CUIs directly; otherwise, normalized-string matching is used as a fallback. The later grouping and canonicalization stages do not depend on CUI assignment; instead, they use the normalized subject and outcome strings produced by the same normalization step. The UMLS-derived entity buckets also contribute to the ILP-based calibration-cluster selection (§\ref{subsec:eval-ilp}), where CUI overlap is one of several weighted relatedness signals. Alternative biomedical linkers and an evaluation benchmark exist~\parencite{}, but the scispaCy/UMLS stack was found sufficient for the needs of this thesis. 